
\subsubsection{2.1 Alignment Methods for Language Models}

Reinforcement learning from human feedback (RLHF) trains language models using reward signals derived from human preferences \citep{ouyang2022training}. Direct Preference Optimization (DPO) provides an alternative formulation that directly optimizes a preference objective without requiring a separate reward model \citep{rafailov2023direct}. Xu et al. (2024) compared PPO and DPO for code generation, reporting that PPO achieved higher pass rates in code competition settings.

\subsubsection{2.2 Reinforcement Learning for Code Generation}

CodeRL \citep{le2022coderl} trains code generation models using binary execution rewards, where code that passes test cases receives reward 1 and code that fails receives reward 0. This creates a reward structure where all non-passing programs---whether they fail to parse, crash at runtime, or produce wrong output---receive identical zero reward. Subsequent work has explored variations including curriculum learning with compiler feedback \citep{dou2024stepcoder} and end-to-end execution feedback \citep{liu2024rlef}.

\subsubsection{2.3 Error Taxonomy in Code Generation}

Prior work has developed taxonomies for classifying code generation errors. The ICSE 2025 study by Wang et al. (2025) provides a three-tier classification based on execution semantics. The LlmFix framework \citep{zhang2024llmfix} extends this to 19 distinct error causes. These taxonomies have been used to characterize errors without stratifying by training method.
